% theme: science_and_technology
\MajorQuestion{科学技術}
\SetRowSpaceQanda

\Instruction{素材と新しい技術に関する次の問いに答えなさい。}
\QQ{石油などを原料として人工的につくられる有機物の素材を\NamedBlank{plastic}という。}{}
\QQ{自然界へ流れ出た\RefBlank{plastic}が細かく砕かれ、粒子になったものを何というか。}{}
\QQ{土壌中の微生物によって分解される\RefBlank{plastic}を何というか。}{}
\QQ{特定のすぐれた性質やはたらきをもつようにつくられた高分子を何というか。}{}
\QQ{コンピュータが学習や判断などを行う技術を、アルファベット2文字で何というか。}{}

\Instruction{発電とエネルギー利用の課題に関する次の問いに答えなさい。}
\QQ{化石燃料を燃やして高温・高圧の水蒸気をつくり、タービンを回す発電方法を何というか。}{}
\QQ{原子核が分裂するときに出るエネルギーを利用して、水蒸気でタービンを回す発電方法を何というか。}{}
\QQ{高い所から流れ落ちる水で水車を回す発電方法を何というか。}{}
\QQ{硫黄酸化物や窒素酸化物などの影響で、強い酸性を示す雨を何というか。}{}
\QQ{大気中の二酸化炭素濃度の上昇などにより、地球の平均気温が上がる現象を何というか。}{}

\Instruction{放射線の性質と単位に関する次の問いに答えなさい。}
\QQ{原子核から出る粒子の流れや電磁波などをまとめて\NamedBlank{radiation}という。}{}
\QQ{\RefBlank{radiation}を出す能力を何というか。}{}
\QQ{\RefBlank{radiation}を出す性質をもつ物質を何というか。}{}
\QQ{物質を透過する性質などをもつ\RefBlank{radiation}を出す能力の大きさを表す単位を何というか。}{}
\QQ{\RefBlank{radiation}が人体に与える影響の大きさを表す単位を何というか。}{}

\Instruction{これからのエネルギー利用と社会に関する次の問いに答えなさい。}
\QQ{太陽光や風力など、使い続けてもなくならないエネルギーを\NamedBlank{renewable_energy}という。}{}
\QQ{発電時に生じる排熱の一部を、給湯や暖房にも利用するしくみを何というか。}{}
\QQ{水素と酸素の化学変化を利用し、化学エネルギーから電気エネルギーを得る装置を何というか。}{}
\QQ{資源の消費を減らし、再利用を進めて、資源を循環させる社会を何というか。}{}
\QQ{生活に必要なものやエネルギーを、現在だけでなく将来の世代も安定して得られる社会を何というか。}{}
